%---------- Inleiding ---------------------------------------------------------

% TODO: Is dit voorstel gebaseerd op een paper van Research Methods die je
% vorig jaar hebt ingediend? Heb je daarbij eventueel samengewerkt met een
% andere student?
% Zo ja, haal dan de tekst hieronder uit commentaar en pas aan.

%\paragraph{Opmerking}

% Dit voorstel is gebaseerd op het onderzoeksvoorstel dat werd geschreven in het
% kader van het vak Research Methods dat ik (vorig/dit) academiejaar heb
% uitgewerkt (met medesturent VOORNAAM NAAM als mede-auteur).
% 

\section{Introduction}%
\label{sec:introduction}

Morden society relies heavily (if not entirely) on available communication networks.
This means emergency reponse relies on this as well. When a major storm destroys a cellular
tower, or flooding cuts power, or large-scale incidents overloads the general public channels,
they become unavailable. Or as we've seen recently during wars and communication lines get
intentionally cut off. The moment they become unavailable, is the exact moment they are needed
the most. In those scenarios first responders, civil protection services and communities need a
fallback that is independent of the existing fixed infrastructure. For those situations this
should be easily and fast deployable solution, without needing any specialist and it
should also be afordable at scale.

A few off-the-grid communication technologies already exist, but they all have their own limitions
in emergency contexts. \textbf{APRS} (Automatic Packet Reporting System) and \textbf{DMR}
(Digital Mobile Radio) both require an amateur radio license and special equiment (specialized hardware
configuration). This makes them hard to deploy in emergency situations. Commercial satellite messengers
and phones (e.g. Garmin inReach or SPOT) are expensive and they require a subscription cost, plus per-message
fees. These limitations make large-scale deployments nearly impossible for small volunteer or
municipal emergency organizations. However, on the other side of the spectrum we also have mesh
radio networks, which are built on \textbf{LoRa} (Long Range) radio technology.
They represent a totally different approach, they are low-power, license-free operating
devices in the \emph{sub-GHz ISM band}, using multi-hop routing which will then extend our coverage
without needing any fixed infrastructure.

To do this we can be looking at \textbf{Meshtastic}.
First, what is Meshtastic?
\begin{itemize}
  \item Open-source firmware and protocol stack for LoRa-capable microcontrollers (mostly ESP32)
  \item Encrypted and decentralized text messaging, locatrion sharing
  \item Self-healing mesh network, multi-hop routing
  \item No need of a SIM card or internet connection, no license
\end{itemize}
Why Meshtastic? On top of the previous mentioned advantages, anecdotes form community deployments and disaster
response suggest promising results. However, no scientific, systematic, empirical research has been done on the
reliability, deployability and evalution of it's performance in emergency situations has yet been published.
This gap in research motivates the following research question:
\textbf{To what extent does Meshtastic, as a decentralized LoRa-based mesh network, offer a reliable and readily deployable emergency communication solution when primary network infrastructure fails?}

\subsection{Problem Domain}%
\label{subsec:problem-domain}
Let's examine and evaluate thje communication requirements that come up when the primary infrastructure fails, and if there are any adequate solutions to fix them.
Three sub-questions guide us:
\begin{itemize}
  \item What gaps exist in emergency contexts when the primary infra becomes unavailable?
  \item How do exisitng alternatives compare to Meshtastioc in terms of accessibility, cost, required expertise, etc. ?
  \item What environmental factors of the Belgian landscape and terrain affect LoRa performances once deployed.
  \item How secure is the Meshtastic encrytion in case of traffic interception?
\end{itemize}

\subsection{Solution Domain}%
\label{subsec:solution-domain}
What are the actual performances of a Meshtastic setup, evaulated through a structured PoC and compare those results.
Three sub-questions emerge here as well:
\begin{itemize}
  \item What range and reliability does a Meshtastic network have across different topologies (direct long-distance
  vs multi-hop network) in a real deployment (East Flanders province)?
  \item How does the Meshtastic security model hold up against potential attacks?
  \item How does it compare with other alternatives in terms of requirements in emergency situations?
\end{itemize}

\subsection{Target audience and scope}%
\label{subsec:target-audience-and-scope}
The main audience are IT professionals responsible for business continuity planning and civil protection communication
in case of emergencies. A secondary audience we could describe are volunteer organizations, municipal safety services and 
network enginess evaluating low-cost off the gris communication options, to be independent of the main societal infrastructure.

The scope of our research is limited by following constraints.
\begin{itemize}
  \item The hardware testing stays limited to Meshtastic due to budget and time constraitns. Alternatives will be evaultaed using litterature only.
  \item The PoC setup and measurtements will onlu be on the EU\_868MHz freqeuncy, this is dictated by \textbf{BIPT} (Belgisch Instituut voor Postdiensten en Telecommunicatie).
\end{itemize}
The geographical focus of the PoC will be the East-Flanders province. Measurements will be taken on the
the direct long-distance link between Ronse and Lochristi (Ghent), this straight connection will cover ~65km.
This is most-definitely unfeasible as a direct LoRa link. That's why we will evalute a multi-hop chain, with nodes in Oudenaarde and maybe Marelbeke (if necessary).
These different setups and measuremnts will help us achieve meaningful results about the reliabiltiy, and limitions of Meshtastic on longer distances.

%---------- Stand van zaken ---------------------------------------------------

\section{Literatuurstudie}%
\label{sec:literatuurstudie}

Hier beschrijf je de \emph{state-of-the-art} rondom je gekozen onderzoeksdomein, d.w.z.\ een inleidende, doorlopende tekst over het onderzoeksdomein van je bachelorproef. Je steunt daarbij heel sterk op de professionele \emph{vakliteratuur}, en niet zozeer op populariserende teksten voor een breed publiek. Wat is de huidige stand van zaken in dit domein, en wat zijn nog eventuele open vragen (die misschien de aanleiding waren tot je onderzoeksvraag!)?

Je mag de titel van deze sectie ook aanpassen (literatuurstudie, stand van zaken, enz.). Zijn er al gelijkaardige onderzoeken gevoerd? Wat concluderen ze? Wat is het verschil met jouw onderzoek?

Verwijs bij elke introductie van een term of bewering over het domein naar de vakliteratuur, bijvoorbeeld~\autocite{Hykes2013}! Denk zeker goed na welke werken je refereert en waarom.

Draag zorg voor correcte literatuurverwijzingen! Een bronvermelding hoort thuis \emph{binnen} de zin waar je je op die bron baseert, dus niet er buiten! Maak meteen een verwijzing als je gebruik maakt van een bron. Doe dit dus \emph{niet} aan het einde van een lange paragraaf. Baseer nooit teveel aansluitende tekst op eenzelfde bron.

Als je informatie over bronnen verzamelt in JabRef, zorg er dan voor dat alle nodige info aanwezig is om de bron terug te vinden (zoals uitvoerig besproken in de lessen Research Methods).

% Voor literatuurverwijzingen zijn er twee belangrijke commando's:
% \autocite{KEY} => (Auteur, jaartal) Gebruik dit als de naam van de auteur
%   geen onderdeel is van de zin.
% \textcite{KEY} => Auteur (jaartal)  Gebruik dit als de auteursnaam wel een
%   functie heeft in de zin (bv. ``Uit onderzoek door Doll & Hill (1954) bleek
%   ...'')

Je mag deze sectie nog verder onderverdelen in subsecties als dit de structuur van de tekst kan verduidelijken.

%---------- Methodologie ------------------------------------------------------
\section{Methodologie}%
\label{sec:methodologie}

Hier beschrijf je hoe je van plan bent het onderzoek te voeren. Welke onderzoekstechniek ga je toepassen om elk van je onderzoeksvragen te beantwoorden? Gebruik je hiervoor literatuurstudie, interviews met belanghebbenden (bv.~voor requirements-analyse), experimenten, simulaties, vergelijkende studie, risico-analyse, PoC, \ldots?

Valt je onderwerp onder één van de typische soorten bachelorproeven die besproken zijn in de lessen Research Methods (bv.\ vergelijkende studie of risico-analyse)? Zorg er dan ook voor dat we duidelijk de verschillende stappen terug vinden die we verwachten in dit soort onderzoek!

Vermijd onderzoekstechnieken die geen objectieve, meetbare resultaten kunnen opleveren. Enquêtes, bijvoorbeeld, zijn voor een bachelorproef informatica meestal \textbf{niet geschikt}. De antwoorden zijn eerder meningen dan feiten en in de praktijk blijkt het ook bijzonder moeilijk om voldoende respondenten te vinden. Studenten die een enquête willen voeren, hebben meestal ook geen goede definitie van de populatie, waardoor ook niet kan aangetoond worden dat eventuele resultaten representatief zijn.

Uit dit onderdeel moet duidelijk naar voor komen dat je bachelorproef ook technisch voldoen\-de diepgang zal bevatten. Het zou niet kloppen als een bachelorproef informatica ook door bv.\ een student marketing zou kunnen uitgevoerd worden.

Je beschrijft ook al welke tools (hardware, software, diensten, \ldots) je denkt hiervoor te gebruiken of te ontwikkelen.

Probeer ook een tijdschatting te maken. Hoe lang zal je met elke fase van je onderzoek bezig zijn en wat zijn de concrete \emph{deliverables} in elke fase?

%---------- Verwachte resultaten ----------------------------------------------
\section{Verwacht resultaat, conclusie}%
\label{sec:verwachte_resultaten}

Hier beschrijf je welke resultaten je verwacht. Als je metingen en simulaties uitvoert, kan je hier al mock-ups maken van de grafieken samen met de verwachte conclusies. Benoem zeker al je assen en de onderdelen van de grafiek die je gaat gebruiken. Dit zorgt ervoor dat je concreet weet welk soort data je moet verzamelen en hoe je die moet meten.

Wat heeft de doelgroep van je onderzoek aan het resultaat? Op welke manier zorgt jouw bachelorproef voor een meerwaarde?

Hier beschrijf je wat je verwacht uit je onderzoek, met de motivatie waarom. Het is \textbf{niet} erg indien uit je onderzoek andere resultaten en conclusies vloeien dan dat je hier beschrijft: het is dan juist interessant om te onderzoeken waarom jouw hypothesen niet overeenkomen met de resultaten.

